\documentclass[11pt]{article}
\usepackage[margin=0.92in]{geometry}
\usepackage[T1]{fontenc}
\usepackage[utf8]{inputenc}
\usepackage{amsmath,amssymb,amsthm}
\usepackage{xcolor}
\usepackage{microtype}
\usepackage[colorlinks=true,linkcolor=blue!55!black,citecolor=blue!55!black,urlcolor=blue!55!black]{hyperref}

\newtheorem{theorem}{Theorem}
\newtheorem{lemma}{Lemma}
\newtheorem{remark}{Remark}
\newcommand{\C}{\mathbb C}
\newcommand{\Sp}{\mathcal S}
\newcommand{\Hp}{H_{\bar f}}
\newcommand{\Hm}{H_f}
\newcommand{\rhoo}{\rho}
\newcommand{\ap}{A_p}
\newcommand{\IMO}{\operatorname{IMO}}

\title{A canonical-weight phase diagram for the Schatten Berger--Coburn problem}
\author{Candidate partial answer to Open Problem 1 in arXiv:2312.06656}
\date{12 August 2026}

\begin{document}
\maketitle

\begin{center}
\fbox{\parbox{0.93\textwidth}{\raggedright\textbf{Status: candidate partial result.}
For the canonical doubling Fock weights $\phi_m(z)=|z|^m$, the Berger--Coburn
question is answered for every $1<p<\infty$ when $m\leq2$, and for all but one
high-$p$ range when $m>2$.  The low endpoint is sharp: the positive
Muckenhoupt window meets an explicit counterexample window exactly, with
logarithmic failure at the boundary.}}
\end{center}

\section{Question and phase diagram}

Let $F^2_m=F^2_{|z|^m}$ and let $H_f:F^2_m\to L^2_{|z|^m}$ be the Hankel
operator with bounded symbol $f$.  Asghari, Hu, and Virtanen~\cite{source}
prove the Berger--Coburn phenomenon for the Hilbert--Schmidt class and ask
whether
\[
 H_f\in\Sp_p\quad\Longleftrightarrow\quad H_{\bar f}\in\Sp_p
 \qquad(1<p<\infty)
\]
on an arbitrary doubling Fock space.  Their Remark~6 identifies
$\rho^{p-2}\in A_p$ as the one missing ingredient in their proof strategy,
where $\rho(z)$ is the radius of a disk of unit $\Delta\phi$-measure.

For $m>2$ set
\[
 p_-(m)=\frac{2m}{m+2},\qquad p_+(m)=\frac{2m}{m-2}.
\]

\begin{theorem}[Canonical-weight phase diagram]\label{thm:phase}
Let $m>0$, $1<p<\infty$, and $f\in L^\infty(\C)$.
\begin{enumerate}
\item If $0<m\leq2$, then $H_f\in\Sp_p$ if and only if
      $H_{\bar f}\in\Sp_p$, and their Schatten norms are comparable.
\item If $m>2$ and $p_-(m)<p<p_+(m)$, the same conclusion holds.
\item If $m>2$ and $1<p\leq p_-(m)$, the bounded Xia symbol
\begin{equation}\label{eq:xia}
 f(z)=\begin{cases}z^{-1},&|z|\geq1,\\0,&|z|<1,
 \end{cases}
\end{equation}
satisfies $H_f\in\Sp_p$ but $H_{\bar f}\notin\Sp_p$.
At $p=p_-(m)$ the obstructing integral diverges logarithmically.
\end{enumerate}
Thus, for $m>2$, only $p\geq p_+(m)$ remains undecided.
\end{theorem}

The two endpoints are H\"older conjugates:
$p_-^{-1}+p_+^{-1}=1$.  This symmetry is suggestive, but it does not by
itself transfer either Schatten membership or nonmembership between the two
ranges.

\section{The exact Muckenhoupt calculation}

The source's canonical-weight lemma states that
\begin{equation}\label{eq:rho}
 \rho(z)\asymp |z|^{1-m/2}\qquad(|z|\gg1).
\end{equation}
Since $\rho$ is positive and comparable on compact sets, for every fixed $p$
the weight $\rho^{p-2}$ is comparable on all of $\C$ to
\[
 w_\gamma(z)=(1+|z|)^\gamma,
 \qquad \gamma=(1-m/2)(p-2).
\]
Membership in $A_p$ is unchanged under two-sided comparability.

\begin{lemma}[Regularized planar powers]\label{lem:power}
For $1<p<\infty$,
\[
 (1+|z|)^\gamma\in A_p(\C)
 \quad\Longleftrightarrow\quad
 -2<\gamma<2(p-1).
\]
\end{lemma}

\begin{proof}
Test the $A_p$ expression on $B(0,R)$ as $R\to\infty$.  If
$\gamma\leq-2$, the average of $w_\gamma$ is respectively
$\asymp R^{-2}\log R$ or $\asymp R^{-2}$, while the dual average supplies
the uncancelled positive power; hence the expression is unbounded.  Applying
the same argument to $w_\gamma^{-1/(p-1)}$ gives necessity of
$\gamma<2(p-1)$, including logarithmic failure at equality.

Conversely, assume the strict inequalities and let $B=B(a,R)$.  If
$R\leq(1+|a|)/2$, both weights are essentially constant on $B$, so their
normalized averages cancel.  Otherwise $|a|<2R+1$ and $B$ lies in a centered
disk of radius $O(R+1)$.  The conditions
$\gamma>-2$ and $-\gamma/(p-1)>-2$ give
\[
 \frac1{|B|}\int_B w_\gamma\lesssim(R+1)^\gamma,
 \qquad
 \frac1{|B|}\int_B w_\gamma^{-1/(p-1)}
 \lesssim(R+1)^{-\gamma/(p-1)}.
\]
Their $A_p$ product is uniformly bounded.
\end{proof}

When $0<m<2$, put $a=1-m/2\in(0,1)$.  Then
$\gamma=a(p-2)>-2$ and $\gamma<2(p-1)$ for every $p>1$; the case $m=2$
has $\gamma=0$.  When $m>2$, put $d=m/2-1>0$, so
$\gamma=-d(p-2)$.  The two strict inequalities in Lemma~\ref{lem:power}
become
\[
 p<2+\frac2d=\frac{2m}{m-2}=p_+(m),
 \qquad
 p>\frac{2d+2}{d+2}=\frac{2m}{m+2}=p_-(m).
\]
Consequently $\rho^{p-2}\in A_p$ precisely in the positive ranges of
Theorem~\ref{thm:phase}.  Remark~6 of the source, using the weighted
Ahlfors--Beurling estimate of Dragi\v cevi\'c~\cite{dragicevic}, then gives
\[
 \|H_{\bar f}\|_{\Sp_p}\lesssim\|H_f\|_{\Sp_p}.
\]
Apply the same statement to $\bar f$ for the reverse implication and norm
comparison.  This proves parts (1) and (2).

\section{The counterexample reaches the lower endpoint}

The source proves much more about the one-sided membership of
\eqref{eq:xia} than its stated $p\leq1$ theorem needs.  Outside a sufficiently
large disk, $f$ is holomorphic on every intrinsic disk
$D(z,\rho(z))$.  Hence its integral distance to analytic functions vanishes
there.  On the remaining compact set it is bounded, and local finiteness of
$\rho^{-2}dA$ gives
\[
 f\in\operatorname{IDA}^{p,2,-2/p}\quad\text{and therefore}\quad
 H_f\in\Sp_p
 \qquad\text{for every }p>0
\]
by Theorem~1.2 of the source.

For $|z|$ large, the direct mean-oscillation calculation in the same proof
gives
\begin{equation}\label{eq:mo}
 MO_{2,1}(f)(z)\asymp\frac{\rho(z)}{|z|^2}.
\end{equation}
Theorem~5.4 of the source says that simultaneous membership of $H_f$ and
$H_{\bar f}$ in $\Sp_p$ is equivalent to
$f\in\IMO^{p,2,-2/p}$.  By \eqref{eq:mo}, its tail is finite exactly when
\begin{align*}
 \int_{|z|>R}\rho(z)^{-2}MO_{2,1}(f)(z)^p\,dA(z)
 &\asymp \int_R^\infty
 r^{(1-m/2)(p-2)-2p+1}\,dr
\end{align*}
is finite.  For $m>2$, the exponent is at least $-1$ exactly when
\[
 p\leq\frac{2m}{m+2}=p_-(m).
\]
At equality it is $-1$.  Thus $f\notin\IMO^{p,2,-2/p}$ throughout this
closed range.  Since $H_f\in\Sp_p$, Theorem~5.4 forces
$H_{\bar f}\notin\Sp_p$, proving part (3).

\section{Scope and upgrade audit}

The result is sharp at $p_-(m)$ but does not settle $p\geq p_+(m)$ for
$m>2$.  Failure of the $A_p$ sufficient condition is not itself a
counterexample.  Schatten nesting cannot propagate the Xia nonmembership
upward, and interpolation supplies no endpoint estimate for a well-defined
symbol-conjugation operator.  Nor does unboundedness of the Beurling transform
automatically produce a bounded symbol with the required IDA/IMO asymmetry.

The natural exterior-holomorphic variants also stop here: a bounded function
holomorphic off a disk has a Laurent expansion whose slowest nonconstant
decay is $1/z$, while $1/z^n$ only weakens the tail for $n>1$.  Eight focused
upgrade attempts, including these high-$p$ routes and a later-literature
search through 12 August 2026, found no defensible extension.  The general
doubling-weight problem likewise remains open.

\begin{thebibliography}{9}
\bibitem{source}
G. Asghari, Z. Hu, and J. A. Virtanen,
\emph{Schatten class Hankel operators on doubling Fock spaces and the
Berger--Coburn phenomenon}, arXiv:2312.06656v2 (2024); J. Math. Anal. Appl.
540 (2024), 128596.

\bibitem{dragicevic}
O. Dragi\v cevi\'c,
\emph{Weighted estimates for powers and the Ahlfors--Beurling operator},
Proc. Amer. Math. Soc. 139 (2011), 2113--2120.
\end{thebibliography}

\end{document}
